\documentclass[./informes.tex]{subfiles}

\begin{document}

\ifSubfilesClassLoaded{
    \chapter{Informes}
}{}

\section{informe-01}

Semana 01: Jueves 19 marzo 2026

En esta primera sesión tuvimos la presentación del curso, donde se presenta el programa del curso, el calendario de lecturas, y las instrucciones para los informes de lectura.

\subsection{Formalidades}

\begin{enumerate}
    \item Debe ser entregado antes de comenzar el desarrollo de un tema.
    \item La extensión máxima es de cuatro páginas, aunque se aprecia que sean más breves de ser posible.
    \item Se considera el cuidado en redacción y ortografía.
    \item Las citas expresas al texto deben ir entrecomilladas y con referencia completa (MLA o APA).
    \item La información extraída de páginas Web también debe ser referida.
\end{enumerate}


\subsection{Contenido}

\begin{enumerate}
    \item El orden en que presenten las ideas es libre, no necesita estar atado al despliegue del texto original.
    \item Se debe consignar de lo que se trata el texto, esto es, su tema fundamental y los objetivos del autor al redactarlo.
    \item Deben incluirse las ideas principales del autor, así como los argumentos que los fundamentan.
    \item Ojo, no es un resumen, el texto ya es conocido y no es necesario que lo repitan todo: se debe distinguir lo principal de lo secundario
    \item Las ideas deben presentarse en forma encadenada, de tal manera que exista una ilación lógica entre ellas.
    \item No se exige toma de posición ni crítica frente al autor y el texto.

\end{enumerate}


\newpage

\end{document}
